\scanpage{38}
complex analytic domain). Here
$X=\mathbb C$ is topologically trivial, which
can thus be expressed by the fact that
its homotopy type is completely
described by that of the topos $\unsure{\mathcal F|\mathfrak F}$ of
systems (17). Indeed, we shall have
\[
(22)\qquad \unsure{\mathcal F|\mathfrak F}\simeq\widehat C
\]
where $C$ (in the general case \deleted{of the}\inserted{where in the} diagram
(19) \deleted{\unsure{with $\pi_1(Y)$ reduced to $\{1\}$}}
\inserted{where $\pi_1(Y)=\{1\}$},
\inserted{$D\Rightarrow\pi_1(\mathcal V^*_{Y,X})\xrightarrow{\sim}\pi_1(X^{\unsure{*|\ell}})$ isomorphism)}
\deleted{\unsure{then}}) is \unsure{defined|defined} as the category whose
set of objects is the set \deleted{\unsure{sum}}\inserted{$\{a,b\}$,}
\deleted{\unsure{union} \illegible{3} \unsure{of the three}}
\deleted{\unsure{arrows of a subcategory}}
\deleted{\unsure{of the sets of objects} of the groupoids, components
\inserted{\unsure{to appear in the} \deleted{\unsure{a}} category}
of (8), and whose arrows are, in addition
to the arrows of these three groupoids, the
arrows $\deleted{\unsure{p}}$ from an object $a$ of $\inserted{A=}\Pi_1\mathcal V^*_{Y,X}$
to an object either of $B=\Pi_1\deleted{\unsure{\mathcal V_{Y,X}|\mathcal V^*_{Y,X}}}\inserted{\unsure{Y}}$, or of
$C=\Pi_1\deleted{\unsure{\mathcal V_{Y,X}|\mathcal V^*_{Y,X}}}\inserted{X^*}$, \unsure{defined|defined} as being
the arrows :
\[
\deleted{\unsure{\varphi}}\quad p_*(a)\longrightarrow b\quad\text{in }B
\quad(\text{resp. }j_*(a)\longrightarrow c\quad\text{in }C).
\]
In the \deleted{\unsure{current}}, present case, the groupoids
$A,B,C$ and the morphisms between them}
% Pale lower-left foliation, in another hand; not the author's page number.
\margin{37}
